% together/seqlock.tex
% mainfile: ../perfbook.tex
% SPDX-License-Identifier: CC-BY-SA-3.0

\section{顺序锁（sequence lock）特技}
\label{sec:together:Sequence-Locking Specials}
%
\epigraph{不会跳舞的姑娘说乐队演奏得不好。}
	 {意第绪谚语}

本节探讨 sequence lock 的一些特殊用法。

\subsection{对决的顺序锁}
\label{sec:together:Dueling Sequence Locks}

经典的 sequence lock 用例使读者能够看到
一小组变量的一致快照，
例如时间保持的校准常数。
这在实践中运行得很好，
因为校准常数很少更新，
而且更新时速度很快。
因此读者几乎不需要重试。

然而，如果更新者在更新期间被延迟，
读者也将被延迟。
这种延迟可能由中断、NMI，
甚至虚拟 CPU 抢占引起。

防止更新者延迟导致读者延迟的一种方法是
维护两组校准常数。
每组依次更新，但更新频率足够高，
使读者可以充分利用任一组。
每组都有自己的 sequence lock（\co{seqlock_t} 结构）。

更新者交替更新两组，
这样被延迟的更新者最多只会延迟其中一组的读者。

每个读者首先尝试访问第一组，
但在需要重试时尝试访问第二组。
如果第二组也迫使重试，
读者从第一组重新开始。
如果更新者卡住了，两组中只有一组会迫使读者重试，
因此读者在尝试访问另一组时就会成功。

\QuickQuiz{
	为什么不是所有 sequence lock 用例
	都以这种方式复制数据？
}\QuickQuizAnswer{
	如果数据太大，这种复制就不切实际，
	例如 \cref{sec:together:Correlated Data Elements}
	中描述的 Schr\"odinger 动物园示例。

	如果延迟被阻止，则不需要这种复制，
	例如当更新者在裸金属硬件上运行时禁用中断
	（即不使用容易发生 vCPU 抢占的虚拟机管理程序）。

	或者，如果读者可以容忍偶尔的延迟，
	则同样不需要复制。
	考虑读写锁的例子，
	写者总是延迟读者，反之亦然。

	然而，如果要复制的数据相当小，
	如果延迟是可能的，
	并且读者无法容忍这些延迟，
	复制数据是一种出色的方法。
}\QuickQuizEnd

\subsection{关联数据元素}
\label{sec:together:Correlated Data Elements}

假设我们有一个 hash table，
我们需要对其中两个或多个元素的关联视图。
这些元素一起更新，
我们不希望看到第一个元素的旧版本
与其他元素的新版本混在一起。
例如，Schr\"odinger 决定将其大家庭连同所有动物
添加到其内存数据库中。
虽然 Schr\"odinger 理解婚姻和离婚不会瞬间发生，
但他也是一个传统主义者。
因此，他绝对不希望其数据库显示
新娘已婚但新郎未婚的情况，反之亦然。
另外，如果你觉得 Schr\"odinger 是传统主义者，
试试和他的一些家庭成员交谈！
换句话说，Schr\"odinger 希望能够对其数据库
进行婚姻一致性遍历。

一种方法是使用 sequence lock
（参见 \cref{sec:defer:Sequence Locks}），
使与婚姻相关的更新在
\co{write_seqlock()} 保护下进行，
而需要婚姻一致性的读取在
\co{read_seqbegin()} / \co{read_seqretry()} 循环内进行。
请注意，sequence lock 并非 RCU 保护的替代品：
sequence lock 防止并发修改，
但仍需要 RCU 来防止并发删除。

当关联元素数量少、读取这些元素的时间短
且更新速率低时，这种方法效果很好。
否则，更新可能发生得太快，读者可能永远无法完成。
虽然 Schr\"odinger 不认为即使他最不理智的亲戚
也会快速到足以使婚姻和离婚成为问题，
但他确实意识到这个问题在其他情况下可能会出现。
避免读者饥饿（starvation）问题的一种方法是让读者在重试次数过多时
使用更新侧原语，
但这会降低性能和可扩展性。
避免\IX{starvation}的另一种方法是使用多个 sequence lock，
在 Schr\"odinger 的情况下，也许每个物种一个。

此外，如果更新侧原语使用过于频繁，
由于锁竞争（lock contention）将导致性能和可扩展性下降。
避免此问题的一种方法是维护每元素 sequence lock，
并在更新婚姻状态时持有双方配偶的锁。
读者可以在任一配偶的锁上进行重试循环，
以获得涉及该对成员的婚姻状态变化的稳定视图。
这避免了由于高结婚和离婚率造成的竞争，
但增加了在单次数据库扫描中
获取所有婚姻状态稳定视图的复杂性。

如果元素分组是明确定义且持久的
（婚姻状态被希望如此），
那么一种方法是向数据元素添加指针，
将给定组的成员链接在一起。
读者然后可以遍历这些指针来访问
与第一个找到的元素在同一组中的所有数据元素。

这种技术在 Linux 内核中被大量使用，
最值得注意的也许是 dcache 子系统~\cite{NeilBrown2015RCUwalk}。
请注意，类似的方案也可能适用于 hazard pointer。

这种方法为成功的读者提供顺序一致性，
每个读者要么看到给定更新的效果，要么看不到，
任何部分更新都会导致读侧重试。
顺序一致性是一种极强的保证，
带来同样强的限制和同样高的开销。
在这种情况下，我们看到读者可能一方面被饥饿，
另一方面可能需要获取更新侧锁。
虽然这在更新不频繁的情况下效果很好，
但它不必要地强制读侧重试，
即使更新不影响重试读者访问的任何数据。
\Cref{sec:together:Correlated Fields} 因此介绍了一种
弱得多的一致性形式，
不仅避免了读者饥饿，还避免了任何形式的读侧重试。
下一节则介绍一种更弱的一致性形式，
可以以更低的读者饥饿概率提供。

\subsection{原子移动}
\label{sec:together:Atomic Move}

假设单个数据元素从一个数据结构移动到另一个，
读者仅查找单个数据结构。
然而，当数据元素移动时，读者决不能同时在两个结构中看到它，
也决不能同时在两个结构中都看不到它。
同时，任何在新位置看到该元素的读者
决不能随后在旧位置看到它。
此外，移动可以通过将旧数据元素的新副本
插入目标位置来实现。

例如，考虑一个支持对读者原子重命名操作的 hash table。
扩展 Schr\"odinger 的动物园，
假设动物的名字发生变化，
例如，Schr\"odinger 传统家庭中的每位新娘
可能将姓氏改为与新郎相同。

但更改名字可能会改变哈希值，
也可能要求新娘的元素从一个哈希链移动到另一个。
上述一致性要求如果读者成功查找新名字，
则该读者随后对旧名字的任何查找必须失败。
类似地，如果读者对旧名字的查找导致查找失败，
则该读者随后对新名字的任何查找必须成功。
简而言之，给定读者不应看到新娘瞬间消失，
也不应该在新名字下查找新娘后又在旧名字下查找到她。

可以使用 \cref{sec:together:Correlated Data Elements} 中描述的
单个全局 sequence lock 来强制执行此一致性保证，
但即使对于不查找当前正在进行名字更改的新娘的读者，
这也可能导致读者饥饿。
此保证也可以通过要求读者获取每哈希链锁来强制执行，
但回顾
\cref{fig:datastruct:Read-Only Hash-Table Performance For Schroedinger's Zoo}
显示，即使对于单插槽系统，
这也会导致较差的性能和可扩展性。

另一种对读者更友好的实现方式是使用 RCU，
并在每个元素上放置一个 sequence lock。
查找给定元素的读者在其对该元素的
整个访问过程中充当 sequence lock 读者。
请注意，这些 sequence lock 操作将对每个读者的查找进行排序。

重命名元素可以大致按以下步骤进行：

\begin{enumerate}
\item	获取保护重命名操作的全局锁。
\item	分配并初始化一个具有新名字的元素副本。
\item	写获取旧名字元素上的 sequence lock，
	这有一个副作用：将此获取与后续的插入操作排序。
	对旧名字的并发查找现在将反复重试。
\item	插入具有新名字的元素副本。
	对新名字的查找现在将成功。
\item	执行 \co{smp_wmb()} 以将前面的插入与后续的删除排序。
\item	删除旧名字的元素。
	对旧名字的并发查找现在将失败。
\item	如有必要（例如锁依赖检查器要求），
	写释放 sequence lock。
\item	释放全局锁。
\end{enumerate}

因此，查找旧名字的读者将不断重试，
直到新名字可用，此时他们的最后一次重试将失败。
对新名字的任何后续查找将成功。
任何成功查找新名字的读者都能保证
对旧名字的任何后续查找将失败，可能需要经过一系列重试。

\QuickQuizSeries{%
\QuickQuizB{
	是否可以在新元素插入之前写获取其 sequence lock，
	而不是在旧元素被删除之前获取旧元素的 sequence lock？
}\QuickQuizAnswerB{
	可以，细节留给读者作为练习。

	术语 \emph{tombstone}（墓碑）有时用于指代
	sequence lock 被获取后具有旧名字的元素。
	类似地，术语 \emph{birthstone}（诞生石）有时用于指代
	sequence lock 仍被持有的具有新名字的元素。
}\QuickQuizEndB

\QuickQuizE{
	是否可以避免使用全局锁？
}\QuickQuizAnswerE{
	可以，一种方法是使用每哈希链锁。
	更新者可以获取与旧元素和新元素对应的锁，
	按地址顺序获取它们。
	在这种情况下，插入和删除操作当然需要
	避免获取和释放这些相同的每哈希链锁。
	如果重命名操作频繁，
	这种复杂性是值得的，
	当然也允许重命名操作并发执行。
}\QuickQuizEndE
}% End of \QuickQuizSeries

当然，也可以使用简单的标志
更高效地实现此过程。
然而，这可以被视为 sequence lock 的简化变体，
它依赖于给定元素的 sequence lock
永远不会被写获取超过一次这一事实。

% @@@ Reference TBD flag-for-deletion section.

\subsection{升级为写者}
\label{sec:together:Upgrade to Writer}

如 \cref{sec:defer:Quasi Reader-Writer Lock} 中所讨论的，
RCU 允许读者升级为写者。
当读者在扫描 RCU 保护的数据结构时
注意到当前元素需要更新时，
此功能非常有用。
当你尝试用 sequence lock 实现这个技巧时会发生什么？

事实证明，这种 sequence lock 技巧实际上在
Linux 内核中被使用，
例如 \path{drivers/infiniband/hw/hfi1/sdma.c}
中的 \co{sdma_flush()} 函数。
效果是注定包含的读者需要重试。
因此，当读者检测到某种需要重试的条件时，
会使用这个技巧。
